\documentclass[12pt]{article}
\usepackage[inner=1in,outer=2.2in,top=1in,bottom=1in,marginparwidth=1.6in,marginparsep=0.15in]{geometry}
\usepackage{setspace}
\usepackage{times}
\usepackage{parskip}
\usepackage{fancyhdr}
\usepackage{xcolor}
\usepackage{marginnote}
\usepackage[colorlinks=true,linkcolor=blue!60!black]{hyperref}
\usepackage{tabularx}
\usepackage{array}
\usepackage{booktabs}
\usepackage{soul}
\sethlcolor{yellow!35}

% Footnote styling
\renewcommand{\thefootnote}{\arabic{footnote}}

% Margin note command for rhetorical moves
\newcommand{\rmove}[1]{\marginnote{\raggedright\footnotesize\color{blue!70!black}\textbf{Rhetorical Move:} #1}[0pt]}

% Header
\pagestyle{fancy}
\fancyhf{}
\rhead{\small Chan}
\rfoot{\small\thepage}
\renewcommand{\headrulewidth}{0pt}
\setlength{\headheight}{14.5pt}
\addtolength{\topmargin}{-2.5pt}

\setlength{\parindent}{0.5in}
\doublespacing

\begin{document}

\begin{center}
{\large \textbf{Differential Inclusion and the Violence of Visibility:}}\\[2pt]
{\large \textbf{The \textit{Fresh Off the Boat} Pilot Episode}}\\[12pt]
Skyler Chan\\
ASA201: Introduction to Asian American Studies\\
Midterm Essay
\end{center}

\vspace{6pt}

\rmove{INTRO: Concrete opening $\rightarrow$ thesis. Jessica's line and the failed Lunchables establish the core contradiction up front.}
The lunchbox is among the most ordinary objects in American childhood---plastic, mass-produced, unremarkable. Yet in the pilot episode of \textit{Fresh Off the Boat}, the lunchbox becomes a site where the boundaries of American belonging are negotiated and enforced. When eleven-year-old Eddie Huang brings homemade Chinese food to his new Orlando middle school and white classmates shout that he is eating worms, his mother Jessica takes him to buy Lunchables. In the grocery store, watching him celebrate fitting into American packaging, she asks a question the show never fully answers: ``You want to fit inside a box? That's so American. Why are you so American?'' But the Lunchables do not protect him. The next day, a classmate steals Eddie's tray and calls him a ``chink.'' \hl{Performing whiteness through consumption does not prevent racialization---it only changes its terms.}\footnote{\textbf{Violence of inclusion} (Week 4 -- Espiritu, Lowe, \& Yoneyama, ``Transpacific Entanglements''): Violence accompanies not only exclusion from the nation but also inclusion into it. Inclusion itself becomes a technology of racial governance.} This is what Espiritu, Lowe, and Yoneyama theorize as the \textit{violence of inclusion}: \hl{Asian Americans are incorporated into visibility on terms that require aestheticizing racist violence into comedy, privatizing structural racism into family narrative, and performing model minority}\footnote{\textbf{Model minority} (Week 2 -- Okihiro, ``Is Yellow Black or White?''): The ideological construction of Asian Americans as hardworking and uncomplaining---used to discipline other communities of color and to justify the withdrawal of social services. The model minority is the ideal subject of differential inclusion: included precisely because they do not threaten the system.} compliance. The pilot reveals the \textit{complexities} of this process---differential inclusion\footnote{\textbf{Differential inclusion} (Week 4 -- Espiritu, Lowe, \& Yoneyama): A mode of governance in which marginalized groups are brought \textit{into} the nation but on terms that maintain their subordination. Unlike exclusion, it grants visibility while requiring compliance with norms that preserve racial hierarchy.} operates not only through the episode's content but through the formal strategies of the sitcom genre itself, and the show's own \textit{contradictions}---moments where its critique escapes its assimilationist frame---suggest what differential inclusion might \textit{imagine} beyond itself.

\rmove{COMPLEXITIES: How diff.\ inclusion works through \textit{form}, not just content. Names 4 formal devices and shows each doing ideological work. This is what the prompt means by ``formal, aesthetic, or rhetorical strategies.''}
The first \textit{complexity} of differential inclusion is that it does not operate through censorship or exclusion---it is far more nuanced than simple silencing. It works through form, through the very aesthetic apparatus that grants visibility in the first place.\footnote{\textbf{Hegemony} (Week 2 -- Lowe, ``Heterogeneity, Hybridity, Multiplicity''): Following Gramsci, Lowe theorizes hegemony as the process through which dominant groups maintain power through \textit{consent}, not force. The sitcom's formal conventions manufacture consent about how to feel about racism.} As a single-camera comedy, the pilot relies on editing rhythm, voiceover narration, and music cues to manage the audience's affective response to racial violence---and this is what makes the process so complicated: \hl{the racism is \textit{shown}, not hidden.} In the cafeteria scene, the show's quick editing moves the audience past Eddie's humiliation before it can register as violence---the comedic pacing tells viewers to process racism as a setback, not an assault. The second complexity is the voiceover: Eddie's adult narration delivers the line ``when you live in a Lunchables world, it's not always easy being homemade Chinese food'' with wry detachment that converts the entire episode's violence into nostalgic humor. The narrator teaches viewers the correct affective distance from racialized exclusion---assuring them he survived, that this is a story being told safely from the other side. Meanwhile, the camera frames Eddie in isolation---an individual setback---rather than connecting his humiliation to the systemic pattern the episode itself reveals: the neighborhood women who marvel that Jessica's English is ``very good,'' the restaurant that needs a white host named Mitch to attract customers.\footnote{\textbf{Aestheticizing commodification} (Week 2 -- Lowe): Mainstream culture reduces Asian American difference to consumable aesthetic markers---food, fashion, ``quirky'' traits---while erasing the material histories of labor migration and racial segregation that actually structure Asian American life.} The sitcom treats these as separate comedic situations rather than evidence of a single system. The third and most fundamental complexity is structural: the twenty-two-minute runtime demands that the lunchbox problem be solved before the credits roll. The genre cannot sustain anger, cannot leave racism unresolved, cannot imagine collective action. It is what Lowe, drawing on Gramsci, would call a hegemonic apparatus: a form that naturalizes assimilation as the only available narrative. \hl{Differential inclusion is complicated precisely because it does not look like oppression---it looks like opportunity, like a seat at the table, like a show on ABC.}

\rmove{CONTRADICTIONS: Where the text undermines its own assimilationist logic. Parallel structure reveals the pattern the show can't see; Lunchables fail on their own terms.}
But differential inclusion is not only complicated---it is also deeply \textit{contradictory}, undermining its own logic at every turn. Remember what differential inclusion promises: comply with the terms of visibility---erase your difference, perform assimilation---and you will be accepted. The pilot exposes this bargain as false from three directions. The first contradiction is the parallel structure between Eddie and Louis. Eddie hides Chinese food behind Lunchables; Louis hires Mitch, a white man, as his front-of-house host because a ``Caucasian face'' will make diners comfortable. Both perform whiteness to enter American space. Differential inclusion frames this as individual choice---Eddie adjusting to a new school, Louis making a smart business decision. But the parallel reveals that it is not a choice at all; it is a \textit{structural demand} imposed at every scale, from a child's lunch table to a father's livelihood.\footnote{\textbf{Racial capitalism} (Week 4 -- Espiritu, Lowe, \& Yoneyama): Capitalism has always been organized through racial hierarchies, extracting differential value from racialized bodies. The Huangs' restaurant exists within a racialized service economy that funnels Asian immigrants into specific labor niches.} This contradicts differential inclusion's presentation of itself as opportunity: \hl{what looks like voluntary assimilation is actually coercion that the system cannot acknowledge as coercion, because doing so would expose that ``inclusion'' requires erasure.} The second contradiction is the food-as-culture trope---Chinese food as ``foreignness,'' Lunchables as ``American.'' Differential inclusion claims to \textit{include} cultural difference, but in practice it demands that difference be eliminated. The very thing supposedly being welcomed---Asian American identity---must be hidden for the welcome to function. This is what Lowe theorizes as aestheticizing commodification: difference is reduced to a consumable binary, and the complex histories of immigration and racial capitalism that actually produced that difference are erased, as though belonging were simply a matter of packaging. And the third and deepest contradiction: Eddie switches to Lunchables and is \textit{still} called a slur. He does exactly what differential inclusion asks---he erases his difference, he performs whiteness through consumption---and the system racializes him anyway. The promised transaction, erase difference and gain acceptance, collapses on its own terms. This is the contradiction at the heart of differential inclusion: \hl{it offers a bargain it structurally cannot honor, because the racial hierarchy it manages does not actually depend on whether the subject complies. Compliance changes the \textit{form} of subordination, not the fact of it.} Differential inclusion is not merely complicated; it is \textit{self-defeating}.

\rmove{IMPLICATIONS + IMAGININGS: The slur scene ruptures the form; Jessica models institutional confrontation; Huang's disavowal \textit{is} the alternative the prompt asks about.}
Yet the pilot is not fully captured by its own machinery, and here the episode's \textit{implications} and \textit{imaginings} emerge. The climactic scene breaks sharply from sitcom tone: the bully tells Eddie ``you're the one at the bottom now'' and calls him a ``chink.'' The Snoop Dogg track cuts out. The cafeteria falls silent. Eddie retaliates physically. For a moment, the comedic editing halts and the show permits the word to land with its full weight---a remarkable choice for an ABC family sitcom.\footnote{\textbf{Collaborative antagonism} (Week 3 -- Chuh, ``Imagine Otherwise: Introduction''): The doubled condition of working together \textit{and} working subversively against. The pilot practices collaborative antagonism: it works within network TV's constraints while occasionally rupturing them.} In the principal's office, the school treats Eddie's response as the problem while ignoring the slur, and Jessica's fury is played not for comedy but for righteous anger: ``That boy called our son a chink. You think that's okay? \ldots We will sue everyone in this school.'' These scenes \textit{imply} what the sitcom form usually suppresses: \hl{that racism is not a misunderstanding between individuals but a structure sustained by institutions.}

The real Eddie Huang's role as the pilot's narrator crystallizes this. Huang's actual voice frames the episode---his name, his story, his perspective are included in the most literal sense. Yet the words he speaks are shaped by ABC's writers' room; \hl{his voice is present but his politics are absent.} The closing voiceover assures viewers ``you don't have to pretend to be someone else in order to belong''---but the entire episode has just proved that you do. Huang publicly called the show an artificial representation, a story diluted for mass consumption, and stopped narrating after Season 1---withdrawing his consent from the very inclusion the network had offered.\footnote{\textbf{Strategic essentialism} (Week 2 -- Lowe via Spivak): Deploying identity categories for political purposes while simultaneously critiquing their limits. Huang used ``Asian American memoir'' to secure the deal, then publicly problematized how that category was deployed.} His disavowal is what differential inclusion \textit{imagines} beyond itself: what Chuh calls collaborative antagonism as a lived practice---using the platform to gain unprecedented visibility, then publicly refusing the terms on which it was granted. \hl{Like Eddie trading his mother's Chinese food for Lunchables, Huang traded his story for a seat at network television's table. Unlike the sitcom's resolution, he found the terms unacceptable and said so.}

\rmove{CONCLUSION: Stakes + the question the essay leaves the reader with.}
The \textit{Fresh Off the Boat} pilot matters not because it is good or bad representation but because it makes visible the mechanism by which representation operates within structures of power.\footnote{\textbf{Multiculturalism's erasure} (Week 3 -- Chuh): Liberal multiculturalism celebrates diversity on the surface while erasing the ongoing structural violence that produces racial difference. The show's existence as a ``diversity win'' for ABC performs this exact function.} The lunchbox scene, the Cattleman's Ranch subplot, the neighborhood microaggressions, and the slur confrontation all stage the same dynamic: Asian Americans are granted entry into American space on the condition that they aestheticize their difference, solve racism privately, and confirm that the system works. This is differential inclusion---visibility on terms that maintain the structures producing inequality. That the show occasionally ruptures its own frame, and that Huang ultimately refused to narrate it, only underscores both how tightly that frame holds and that it need not hold forever. \hl{The question for Asian American cultural politics is not whether to seek representation but whether to accept the terms on which it is offered---and what it looks like to refuse them.}

%% ============================================================
%% STUDY MATERIAL BELOW — NOT PART OF THE ESSAY
%% ============================================================

\newpage
\singlespacing
\newgeometry{margin=0.7in}
\pagestyle{empty}

\begin{center}
{\Large \textbf{CHEAT SHEET: Prompt Evaluation \& Concept Map}}\\[6pt]
{\small For study only --- not part of the handwritten essay}
\end{center}

\vspace{6pt}

\subsection*{How the Essay Meets the Prompt}

\renewcommand{\arraystretch}{1.4}
\begin{tabularx}{\textwidth}{|>{\raggedright\arraybackslash\bfseries}p{3.8cm}|X|>{\centering\arraybackslash}p{1.2cm}|}
\hline
\textbf{Prompt Requirement} & \textbf{Where / How the Essay Does It} & \textbf{Grade} \\
\hline
Key term from list & Differential inclusion (primary), with hegemony, collaborative antagonism, strategic essentialism, model minority, racial capitalism woven in & Strong \\
\hline
Cultural text analyzed & \textit{Fresh Off the Boat} pilot episode (S1E1) + Huang's narration/disavowal as paratext & Strong \\
\hline
Complexities of key term & Diff.\ inclusion operates through \textit{form} (genre constraints), not just content; it works through consent, not exclusion & Strong \\
\hline
Contradictions of key term & Closing VO says ``you don't have to pretend'' but the episode proves you do; Lunchables don't protect Eddie; Huang's voice is included but his politics are absent & Strong \\
\hline
Implications of key term & Representation $\neq$ liberation; visibility has costs; the question is not whether to seek representation but on what terms & Strong \\
\hline
Imaginings of key term & Huang's disavowal as collaborative antagonism --- using visibility then refusing its terms. Conclusion now explicitly frames this as what diff.\ inclusion ``imagines beyond itself.'' & Strong \\
\hline
Formal/aesthetic/rhetorical strategies & Editing rhythm, voiceover narration, parallel structure, tonal rupture/silence, food as cultural signifier, 22-min runtime, Huang as narrator & Strong \\
\hline
Thesis-driven & Clear 3-part thesis in intro; every paragraph advances it & Strong \\
\hline
Bulk on cultural text (not theory) & $\sim$85\% cultural text analysis, $\sim$15\% theoretical framing & Strong \\
\hline
Originality & Parallel structure as critical juxtaposition; Huang's disavowal; voiceover contradiction; Lunchables fail to protect Eddie & Strong \\
\hline
4--5 handwritten pages & Body text is $\sim$1,100 words. This is a study document---condense as needed when writing by hand. The margin notes label which sections to prioritize. & N/A \\
\hline
\end{tabularx}

\vspace{12pt}
\subsection*{Differential Inclusion $\times$ Formal Strategies: The Core Map}

{\small The prompt asks: what does the text show about the key term's \textbf{complexities}, \textbf{contradictions}, \textbf{implications}, and \textbf{imaginings}? This table maps each formal strategy to all four.}

\vspace{6pt}

\renewcommand{\arraystretch}{1.3}
\footnotesize
\begin{tabularx}{\textwidth}{|>{\raggedright\arraybackslash\bfseries}p{1.8cm}|X|X|X|X|}
\hline
\textbf{Strategy} & \textbf{Complexity} \newline {\scriptsize (how diff.\ inclusion is more nuanced than it seems)} & \textbf{Contradiction} \newline {\scriptsize (where it undermines itself)} & \textbf{Implication} \newline {\scriptsize (what this means politically)} & \textbf{Imagining} \newline {\scriptsize (what alternatives emerge)} \\
\hline

Editing rhythm / comedic pacing &
Diff.\ inclusion doesn't require censorship---it uses \textit{tempo}. Violence is shown but the pace prevents it from registering as violence. Consent is manufactured through speed, not suppression. &
The audience sees the racism but is trained not to feel it. The form includes the content of exclusion while excluding the affect. &
If racism can be \textit{shown} without being \textit{felt}, then visibility alone cannot produce political consciousness. Representation without affect is inert. &
The slur scene breaks the pace---proving the form \textit{can} slow down. The alternative is editing that lets violence land. \\
\hline

Voiceover narration &
The narrator is both inside and outside the story. His retrospective tone creates distance that makes racial violence safe to consume---``I survived, so you can laugh.'' &
Huang's scripted VO says ``you don't have to pretend to be someone else to belong''---but the episode just proved you do. The narrator contradicts his own story. &
If the narrator must assure the audience that everything turned out fine, diff.\ inclusion requires the \textit{past tense}---racism is always ``back then,'' never structural and ongoing. &
Huang stops narrating after S1---refusing to keep providing the reassuring voice. Withdrawing the VO is an act of refusing to make inclusion comfortable. \\
\hline

Parallel structure (lunchbox $\leftrightarrow$ restaurant) &
Diff.\ inclusion operates at every scale simultaneously: the child's lunch table and the father's business both require performing whiteness for entry. It's not isolated---it's systemic. &
The show presents both as solutions (Lunchables work, Mitch works) but reading them together reveals both as the \textit{same erasure} at different scales. The show can't see its own pattern. &
If both a child and his father must independently perform the same racial transaction, then diff.\ inclusion is not a personal challenge but a structural condition reproduced across generations. &
Jessica's line ``Why are you so American?'' is the show's own counter-voice---it names the cost even as the plot endorses it. \\
\hline

Tonal rupture / silence (slur scene) &
Diff.\ inclusion is never total. The sitcom form \textit{almost} contains all critique, but the slur scene proves the form has seams---moments where the political breaks through. &
The show permits the word ``chink'' to land with full weight on ABC---but then restores optimism with Louis's speech. The rupture is real but temporary. &
That the principal punishes Eddie (not the bully) reveals institutional racism. Diff.\ inclusion operates through institutions that treat the response to racism as worse than racism itself. &
Jessica's fury in the principal's office models an alternative: refusing to accept institutional framing, threatening to sue, naming the slur as the real problem. \\
\hline

Food as cultural signifier &
Diff.\ inclusion commodifies difference: Chinese food = ``exotic otherness'' and Lunchables = ``American belonging.'' The binary is aesthetic, not material---it erases \textit{why} Chinese food is stigmatized. &
Eddie switches to Lunchables and is \textit{still} called a slur. The promised transaction (erase difference $\rightarrow$ gain acceptance) fails on its own terms. Assimilation doesn't deliver what it promises. &
If the marketplace (buying Lunchables) is the site of assimilation, then diff.\ inclusion operates through racial capitalism: belonging is purchased, and the purchase erases the history that produced the stigma. &
Jessica's homemade Chinese food represents what diff.\ inclusion demands you give up. Its return at the end (``it's what makes you special'') gestures toward valuing what the system stigmatizes. \\
\hline

22-minute episodic structure &
Diff.\ inclusion works through genre: the sitcom's runtime \textit{structurally requires} resolution. Racism must be ``solved'' before credits roll. The form forecloses sustained political critique. &
The episode ``resolves'' the lunchbox problem, but the slur scene happens \textit{after} the resolution---meaning the real violence occurs after the supposed fix. Resolution is false. &
If the genre cannot sustain anger or imagine collective action, then the sitcom form is itself a hegemonic apparatus---it naturalizes assimilation as the only available narrative. &
A different genre (documentary, drama, limited series) could hold the unresolved tension. The sitcom's limits point toward what other forms might make possible. \\
\hline

Huang as narrator + disavowal &
Diff.\ inclusion is not just about fictional characters---it structures the \textit{production} of culture. Huang's inclusion (name, voice, story) on ABC is the real-world version of Eddie's inclusion at the lunch table. &
The show needs Huang's authenticity to exist as ``the first Asian American family sitcom''---but cannot accommodate the anger that \textit{is} his authenticity. It needs him and must neutralize him simultaneously. &
If the creator of the cultural text disavows his own product, then the terms of inclusion are visible: you may be present, but you cannot control your own representation. The apparatus owns your story. &
Huang's disavowal \textit{is} the alternative. Collaborative antagonism: use the platform to gain visibility, then publicly refuse the terms. Keep the category ``Asian American'' in productive crisis rather than letting institutions settle its meaning. \\
\hline

\end{tabularx}

\vspace{12pt}
\subsection*{Key Lines to Remember}

\begin{itemize}\setlength{\itemsep}{2pt}
\item \textbf{Jessica (grocery store):} ``You want to fit inside a box? That's so American. Why are you so American?''
\item \textbf{Bully:} ``Get used to it. You're the one at the bottom now \ldots It's my turn, chink!''
\item \textbf{Jessica (principal's office):} ``That boy called our son a chink. You think that's okay? \ldots We will sue everyone in this school. So fast, it'll make your head spin. Hey, it's the American way, right?''
\item \textbf{Closing VO:} ``When you live in a Lunchables world, it's not always easy being homemade Chinese food. But it's also what makes you special \ldots you don't have to pretend to be someone else in order to belong.''
\item \textbf{Louis (wild west speech):} ``I opened a wild west restaurant because this is the wild west---a lawless land for only the bravest of families!''
\end{itemize}

\vspace{12pt}
\subsection*{Quick Reference: The 3s of Differential Inclusion}

\renewcommand{\arraystretch}{1.4}
\normalsize

\textbf{Three Core Mechanisms of Differential Inclusion} {\small (from thesis --- what it \textit{requires})}

\begin{enumerate}\setlength{\itemsep}{2pt}
\item \textbf{Aestheticizing racist violence into comedy} --- the sitcom's editing, pacing, and voiceover convert racial abuse into comedic beats, making violence consumable rather than confrontational.
\item \textbf{Privatizing structural racism into family narrative} --- instead of institutions addressing racism, the Huangs solve the problem internally (buying Lunchables, hiring Mitch). Racism becomes a family challenge, not a systemic condition.
\item \textbf{Performing model minority compliance} --- the Huangs are hardworking, uncomplaining, and ask nothing of institutions. They adapt rather than organize --- exactly the subject differential inclusion rewards.
\end{enumerate}

\vspace{8pt}

\textbf{Three Main Complexities} {\small (how diff.\ inclusion is more nuanced than it seems)}

\begin{enumerate}\setlength{\itemsep}{2pt}
\item \textbf{It works through consent, not censorship} --- the racism \textit{is} shown, but editing rhythm and voiceover manage the audience's affect so it registers as comedy, not violence. Hegemony through form.
\item \textbf{It operates at every scale simultaneously} --- Eddie hides Chinese food behind Lunchables; Louis hides Chinese ownership behind ``Cattleman Mitch.'' The same racial transaction is demanded of child and father alike, revealing a systemic pattern, not an individual setback.
\item \textbf{The genre itself is the apparatus} --- the 22-minute runtime structurally requires resolution. The sitcom cannot sustain anger, leave racism unresolved, or imagine collective action. Assimilation is the only narrative the form permits.
\end{enumerate}

\vspace{8pt}

\textbf{Three Main Contradictions} {\small (where diff.\ inclusion undermines itself)}

\begin{enumerate}\setlength{\itemsep}{2pt}
\item \textbf{Lunchables fail --- compliance doesn't deliver acceptance.} Diff.\ inclusion promises: erase your difference and you'll be included. Eddie does exactly that --- and is \textit{still} called a slur. The bargain collapses on its own terms. The racial hierarchy doesn't depend on whether the subject complies; compliance changes the \textit{form} of subordination, not the fact of it.
\item \textbf{The voiceover contradicts the story --- inclusion must deny its own evidence.} The closing narration says ``you don't have to pretend to be someone else in order to belong,'' but the episode just proved that you do. Diff.\ inclusion requires this denial: the narrator must assure the audience the system works, even when the plot has demonstrated that it doesn't.
\item \textbf{The show needs Huang but must neutralize him --- inclusion requires erasing the thing it includes.} Diff.\ inclusion claims to welcome Asian American stories, but the series cannot accommodate the anger that \textit{is} Huang's story. It needs his authenticity to market itself and must strip that authenticity to function. This is the contradiction at the core of differential inclusion: it offers inclusion on terms that hollow out what is being included.
\end{enumerate}

\vspace{8pt}

\textbf{Three Main Implications} {\small (what this means politically)}

\begin{enumerate}\setlength{\itemsep}{2pt}
\item \textbf{Visibility $\neq$ liberation} --- representation that requires erasing structural critique is not progress; it is managed inclusion that sustains the system granting it.
\item \textbf{Institutions protect the structure, not the victim} --- the principal punishes Eddie's retaliation, not the slur. The school treats the \textit{response} to racism as worse than racism itself.
\item \textbf{The marketplace replaces politics} --- Jessica buys Lunchables; belonging is purchased through consumption. Differential inclusion routes the path to acceptance through racial capitalism rather than institutional change.
\end{enumerate}

\vspace{8pt}

\textbf{Three Main Imaginings} {\small (what alternatives emerge from the text)}

\begin{enumerate}\setlength{\itemsep}{2pt}
\item \textbf{Huang's disavowal as collaborative antagonism} --- he used the platform to gain unprecedented Asian American visibility on network TV, then publicly refused the terms on which it was granted. Using inclusion strategically while refusing to let it settle.
\item \textbf{Jessica's fury in the principal's office} --- she refuses the institutional framing, names the slur as the real problem, and threatens to sue. She models confrontation with the system rather than private adaptation.
\item \textbf{The slur scene ruptures the form} --- the music cuts, the cafeteria goes silent, comedy halts. The show proves its own form \textit{can} break open to let violence register as violence. The alternative exists within the text itself.
\end{enumerate}

\vspace{14pt}
\subsection*{Fill-in-the-Blank: ``\_\_\_\_ DI because \_\_\_\_''}

{\small Use these as sentence starters when writing by hand. Each one is a ready-made analytical move.}

\vspace{6pt}
\renewcommand{\arraystretch}{1.5}
\begin{tabularx}{\textwidth}{|>{\centering\arraybackslash}p{1.3cm}|>{\raggedright\arraybackslash\bfseries}p{1.8cm}|X|}
\hline
\textbf{Type} & \textbf{Strategy} & \textbf{``\_\_\_\_ DI because\ldots''} \\
\hline
\multicolumn{3}{|l|}{\textbf{COMPLICATES DI because\ldots}} \\
\hline
{\scriptsize FORMAL} & Editing rhythm &
\ldots DI does not censor racism---it \textit{shows} it, but the comedic pacing moves the audience past it before it registers as violence. Power works through consent, not suppression. \\
\hline
{\scriptsize RHETOR.} & Voiceover narration &
\ldots the adult narrator converts racial violence into nostalgic humor, teaching viewers the ``correct'' affective distance. DI manages \textit{how you feel} about racism, not whether you see it. \\
\hline
{\scriptsize STRUCT.} & 22-min runtime &
\ldots the genre structurally requires resolution. DI doesn't need to argue against anger---the form simply cannot hold it. The apparatus does the ideological work automatically. \\
\hline
{\scriptsize STRUCT.} & Parallel structure &
\ldots the same erasure is demanded at every scale (child's lunch table, father's restaurant), revealing DI as systemic, not individual---but the show frames both as personal choices. \\
\hline
\multicolumn{3}{|l|}{\textbf{CONTRADICTS DI because\ldots}} \\
\hline
{\scriptsize NARR.} & Lunchables fail &
\ldots Eddie complies---erases his difference, performs whiteness---and is \textit{still} called a slur. DI promises acceptance in exchange for compliance, but the hierarchy doesn't depend on compliance. The bargain is structurally false. \\
\hline
{\scriptsize RHETOR.} & VO vs.\ plot &
\ldots the closing narration says ``you don't have to pretend to be someone else to belong,'' but the episode just proved you do. DI must deny its own evidence to function. \\
\hline
{\scriptsize AESTH.} & Food trope &
\ldots DI claims to \textit{include} difference, but demands difference be erased. The thing being ``welcomed'' (Asian American identity) must be hidden for the welcome to work. Inclusion requires the elimination of what it includes. \\
\hline
{\scriptsize RHETOR.} & Huang as narrator &
\ldots the show needs Huang's authenticity to market itself as Asian American representation, but cannot accommodate the anger that \textit{is} his authenticity. It includes him and neutralizes him simultaneously. \\
\hline
\multicolumn{3}{|l|}{\textbf{IMPLIES (about DI) that\ldots}} \\
\hline
{\scriptsize NARR.} & Principal's office &
\ldots institutions protect the racial structure, not the victim. The school punishes Eddie's response, not the slur. DI is maintained by institutions that treat the reaction to racism as worse than racism itself. \\
\hline
{\scriptsize AESTH.} & Marketplace / Lunchables &
\ldots belonging is routed through consumption (buying Lunchables), not institutional change. DI operates through racial capitalism: acceptance is purchased, and the purchase erases the history that produced the stigma. \\
\hline
{\scriptsize STRUCT.} & Show on ABC &
\ldots representation $\neq$ liberation. A show on ABC is not progress if it requires erasing structural critique. DI makes visibility itself a tool of governance, not a path out of it. \\
\hline
\multicolumn{3}{|l|}{\textbf{IMAGINES (beyond DI) that\ldots}} \\
\hline
{\scriptsize RHETOR.} & Huang's disavowal &
\ldots you can use the platform \textit{and} refuse its terms. Collaborative antagonism as lived practice: gain visibility, then publicly reject the conditions. Keep ``Asian American'' in productive crisis rather than letting institutions settle its meaning. \\
\hline
{\scriptsize NARR.} & Jessica's fury &
\ldots confrontation with institutions is possible. She names the slur, threatens to sue, refuses the principal's framing. She models demanding accountability from the system rather than solving racism privately. \\
\hline
{\scriptsize FORMAL} & Slur scene / tonal rupture &
\ldots the sitcom form \textit{can} break open. The music cuts, the comedy halts, the word lands with its full weight. The alternative to managed inclusion exists within the text itself---the form just usually won't permit it. \\
\hline
\end{tabularx}

\vspace{6pt}
{\footnotesize \textbf{Key:} FORMAL = how the show is \textit{made} (editing, sound, camera) | RHETOR. = how the text \textit{persuades} (narration, framing, audience positioning) | STRUCT. = genre/narrative \textit{architecture} (runtime, parallel plots, platform) | AESTH. = visual/cultural \textit{signifiers} (food, consumption, objects) | NARR. = \textit{plot events} that reveal DI's logic}

\end{document}
